\uuid{g2fK}
\titre{Éléments simples}
\auteur{bodin ; hanani ; mzari}
\organisation{exo7}
\licence{CC BY-NC-SA 4.0}
\datecreate{2020-01-20}
\niveau{L1}
\module{Algèbre}
\chapitre{Polynôme, fraction rationnelle}
\sousChapitre{Fraction rationnelle}
\difficulte{3}
\theme{Fractions rationnelles, éléments simples}
\exo7qcmid{lille-1.polynomes.029}

\contenu{
\begin{question}
\qclassification{105.04}
\qtags{motcle=éléments simples}
Soit $F(X) = \dfrac{1}{(X^2+1)X^3}$.
On écrit
$$F(X) = \frac{a}{X} + \frac{b}{X^2} +  \frac{c}{X^3}
+ \frac{dX+e}{X^2+1}.$$
Quelles sont les affirmations vraies ?
\begin{answers}
    \good{$c=1$}

    \bad{$b=1$}

    \bad{$a=1$}

    \good{$e=0$}
\end{answers}
\begin{explanations}
On profite que $F$ est impaire pour déduire $b=0$, $e=0$.
$F(X) = \dfrac{1}{(X^2+1)X^3} = \frac{-1}{X}  +  \frac{1}{X^3}
+ \frac{X}{X^2+1}.$
\end{explanations}
\end{question}
}
